%%%%%%%%%%%%%%%%%%%%%%%%%%%%%%%%%%%%%%%%%%%%%%%%%%%%%%%%%%%%%%%%%%
%% Toto je inkludovaný LaTeXový soubor `ca_lisen.tex'.          %%
%% ------------------------------------------------------------ %%
%% Verze 2022-06-05                                             %%
%% Vyvíjí & udržuje <stanislav.hledik@physics.slu.cz>           %%
%%%%%%%%%%%%%%%%%%%%%%%%%%%%%%%%%%%%%%%%%%%%%%%%%%%%%%%%%%%%%%%%%%

\chapter{Výčtová prostředí}\label{lisen}%%%%%%%%%%%%%%%%%%%%%%%%%%%%%%%%%%%

Stejně jako standardní \LaTeX ové\INDEX{publikační systém!\LaTeX{}} třídy,
poskytuje i~dokumentní třída balíčku \pkg{\ipotname} výčtová prostředí
\pgm{itemize} (nečíslovaný výčet), \pgm{enumerate} (číslovaný výčet) a
\pgm{description} (slovníkový výčet)~\cite{Ryb:2005:LaTeX:3,%
  Kop-Dal:2004:LaTeXPodrPruv:}. Navíc poskytuje jejich \uv{hvězdičkové
  varianty} \pgm{itemize*}, \pgm{enumerate*} a \pgm{description*}.

Obyčejná (nehvězdičková) varianta nečíslovaného a číslovaného výčtu a obě
varianty slovníkového výčtu používají \emph{zavěšenou indentaci}, kdy symbol
položky u~nečíslovaného výčtu, číslo položky u~číslovaného výčtu nebo heslo
slovníkového výčtu vyčnívá doleva a další řádky textu položky jsou zarovnány
na začátek prvního řádku textu nečíslované nebo číslované položky resp. na
odstavcovou zarážku u~slovníkového výčtu.

Výčtová prostředí lze do sebe vnořovat. Následující příklad demonstruje takové
vnoření a zároveň odkazuje na příklady různých typů výčtových prostředí
roztroušených v tomto manuálu:

\begin{description}
\item[Nečíslovaný výčet:]\label{edesciitem} \mbox{}\par
  \begin{itemize}
  \item\label{iitema} Příklad nečíslovaného výčtu najdete v Sekci~\ref{sono}
    na stránce~\pageref{sono}.
  \item\label{iitemb} Další příklad nečíslovaného výčtu najdete na konci
    Kapitoly~\ref{princtwo} za rovnicí~(\ref{mrift}).
  \item\label{iitemc} Tento nečíslovaný výčet byl vytvořen pomocí prostředí
    \pgm{itemize} vnořeného do vnějšího slovníkového výčtu vytvořeného pomocí
    prostředí \pgm{description}.
  \end{itemize}
\item[Číslovaný výčet:]\label{edescienum} \mbox{}\par
  \begin{enumerate}
  \item\label{ienuma} Příklad číslovaného výčtu najdete v~Úvodu na
    stránce~\pageref{easyuse}.
  \item\label{ienumb} Další příklad číslovaného výčtu najdete
    v~Sekci~\ref{instsepar} na stránce~\pageref{instsepar}.
  \item\label{ienumc} A ještě další příklady číslovaného výčtu najdete
    v~Kapitole~\ref{CAPczm} před začátkem Sekce~\ref{SECpafpo} a za
    rovnicí~(\ref{HOMmtxM}), v~Kapitole~\ref{SpinLev} před začátkem
    Sekce~\ref{LeviStab}.
  \item\label{ienumd} Tento číslovaný výčet byl vytvořen pomocí prostředí
    \pgm{enumerate} vnořeného do vnějšího slovníkového výčtu vytvořeného
    pomocí prostředí \pgm{description}.
  \end{enumerate}
\item[Slovníkový výčet:]\label{edescidesc} \mbox{}\par
  \begin{description}
  \item[Příklady slovníkového výčtu:]\label{idesca} najdete je
    v~Dodatku~\ref{options} na stránce~\pageref{options} a
    v~Dodatku~\ref{packages} stránce~\pageref{packages} (ačkoli v těchto
    případech je kvůli speciálnímu pořadavku na font hesla použita upravená
    forma příkazu \pgm{\textbackslash item}).
  \item[Tento výčet:]\label{idescb} byl vytvořen pomocí prostředí
    \pgm{description} vnořeného do vnějšího slovníkového výčtu vytvořeného
    rovněž pomocí prostředí \pgm{description}.
  \end{description}
\item[Tento vnější slovníkový výčet:]\label{edesc} vytvořen prostředím
  \pgm{description}.
\end{description}

\noindent
Hvězdičková varianta nečíslovaného a číslovaného výčtu používá indentaci
stejnou jako u normálního odstavce, zatímco u slovníkového výčtu ponechá
původní indentaci a místo toho změní tučný font hesla na kurzívu (italiku):

\begin{description*}
\item[Nečíslovaný výčet:]\label{staredesciitem} \mbox{}\par
  \begin{itemize*}
  \item\label{stariitema} Příklad nečíslovaného výčtu najdete v
    Sekci~\ref{sono} na stránce~\pageref{sono}.
  \item\label{stariitemb} Další příklad nečíslovaného výčtu najdete na konci
    Kapitoly~\ref{princtwo} za rovnicí~(\ref{mrift}).
  \item\label{stariitemc} Tento nečíslovaný výčet byl vytvořen pomocí
    prostředí \pgm{itemize} vnořeného do vnějšího slovníkového výčtu
    vytvořeného pomocí prostředí \pgm{description}.
  \end{itemize*}
\item[Číslovaný výčet:]\label{staredescienum} \mbox{}\par
  \begin{enumerate*}
  \item\label{starienuma} Příklad číslovaného výčtu najdete v~Úvodu na
    stránce~\pageref{easyuse}.
  \item\label{starienumb} Další příklad číslovaného výčtu najdete
    v~Sekci~\ref{instsepar} na stránce~\pageref{instsepar}.
  \item\label{starienumc} A ještě další příklady číslovaného výčtu najdete
    v~Kapitole~\ref{CAPczm} před začátkem Sekce~\ref{SECpafpo} a za
    rovnicí~(\ref{HOMmtxM}), v~Kapitole~\ref{SpinLev} před začátkem
    Sekce~\ref{LeviStab}.
  \item\label{starienumd} Tento číslovaný výčet byl vytvořen pomocí prostředí
    \pgm{enumerate} vnořeného do vnějšího slovníkového výčtu vytvořeného
    pomocí prostředí \pgm{description}.
  \end{enumerate*}
\item[Slovníkový výčet:]\label{staredescidesc} \mbox{}\par
  \begin{description*}
  \item[Příklady slovníkového výčtu:]\label{staridesca} najdete je
    v~Dodatku~\ref{options} na stránce~\pageref{options} a
    v~Dodatku~\ref{packages} stránce~\pageref{packages} (ačkoli v těchto
    případech je kvůli speciálnímu pořadavku na font hesla použita upravená
    forma příkazu \pgm{\textbackslash item}).
  \item[Tento výčet:]\label{staridescb} byl vytvořen pomocí prostředí
    \pgm{description} vnořeného do vnějšího slovníkového výčtu vytvořeného
    rovněž pomocí prostředí \pgm{description}.
  \end{description*}
\item[Tento vnější slovníkový výčet:]\label{staredesc} vytvořen prostředím
  \pgm{description}.
\end{description*}

Číslovaný výčet je vhodný pro kratší seznamy (do 9~položek), neboť pro
dvojciferná pořadová čísla dojde k~nehezkému posunu začátku prvního řádku
položky.

\endinput

Hvězdičkové varianty výčtů \pgm{itemize*}, \pgm{enumerate*} používají
indentaci podobnou jako u~obyčejných odstavců, tj. symbol položky
u~nečíslovaného výčtu a pořadové číslo u~číslovaného výčtu jsou odsazeny
o~odstavcovou zarážku doprava a text na ně plynule navazuje, přičemž další
řádky položky začínají od původního levého okraje. Vnořené výčty jsou celkově
posouvány o~odstavcovou zarážku doprava. Jako příklad uveďme výčet
v~Dodatku~\ref{szzk} na stránce~\pageref{szzk}. Tento způsob je vhodný pro
delší seznamy (více než 9~položek) s~rozsáhlejšími položkami.  Obě varianty
není vhodné ve vnořených seznamech kombinovat.

Hvězdičková varianta slovníkového výčtu \pgm{description*} se co do indentace
chová stejně jako obyčejná, liší se kurzivním zdůrazněním hesla místo tučného.

\endinput

%%
%% End of file `ca_lisen.tex'.


Vnořené výčty jsou celkově posouvány o~odstavcovou
zarážku doprava.